\scanpage{28}
\mspage{27}
On doit donc poser

\uline{A\deleted{\unsure{C}}\unsure{t|E|C}7}\quad Si $X,Y,Z\in\unsure{M|\mathcal{M}}$ sont tels que
\[
 X\overset{\circ}{\ll}Y\overset{\circ}{\ll}Z\qquad\text{alors}\qquad X\overset{\circ}{\ll}Z
\]
i.e. la relation $\overset{\circ}{\ll}$ est \uline{transitive}.

Ceci \unsure{assure|conserve} donc $\Rightarrow$ dans la prop.
En sens inverse, S'il \unsure{existe|exist} $f$ comme dans (1.42), soit $X\in\widetilde G$, \deleted{alors $\varphi_G(\unsure{X|x|/}$} \inserted{\deleted{\unsure{on|ce}}} de sorte que $\varphi_G(\unsure{X|x})=X$ (tautologique\deleted{\unsure{?|!}}), moyennant (6), \deleted{on a alors} soit $f(X)=Y$, on a alors $Y=\varphi_F(X)$ i.e. $Y$ est le plus petit élément de $\widetilde F$ \deleted{\unsure{contenant} $X$}. \deleted{\unsure{Donc}} tel que $X\ll Y$. Cela \deleted{signifie} implique donc que pour tt $X\in\deleted{\widetilde F}\widetilde G$, \unsure{existe|exist} $Y\in\deleted{\unsure{\widetilde G|\widetilde C}}\widetilde F$ telle que $X\ll Y$, i.e. $\deleted{\unsure{\widetilde G|\widetilde F}}$ (C4) $\widetilde G\ll\widetilde F$\unsure{..|:}, qed. On n'a pas \unsure{utilisé|utilisé-|utilis-} en fait que $f$ soit croissante, donc

\uline{Corollaire \unsure{à|2} la prop.}\quad \deleted{La condition,} \inserted{Pour} que $F\ll G$ i.e. $\mathrm{Mulom}(F)\ll\mathrm{Mulomb}(G)$. il s. que $\mathrm{Omb}\,F\subset\mathrm{Omb}\,G$,
\margin{Ce ne sera sans doute plus vrai, en \unsure{passant} aux \uline{petites} multifigures}
et que \deleted{\unsure{toute classe}} la relation d'équivalence \inserted{$\unsure{\mathcal R_F|R_F|Q_F}$} sur $\mathrm{Omb}\,F$ définie par $\varphi_F$ soit plus fine que la relation d'équiv. induite \unsure{par|sur} $\mathrm{Omb}\,G$.
